\documentclass[11pt]{article}

\usepackage[margin=1in]{geometry}
\usepackage{amsmath}
\usepackage{booktabs}
\usepackage{hyperref}

\title{Predicting Short-Term Recovery After Large Stock-Price Crashes}
\author{Studio R42}
\date{\today}

\begin{document}
\maketitle

\section{Motivation}

After a large single-day decline, some stocks rebound within weeks while others
keep falling. We ask a single question: at the moment a stock crashes, using only
information available at that time, can we predict \emph{which} stocks recover, by
how much, and how quickly? A crash is defined as a daily adjusted return of
$-10\%$ or worse, and every qualifying day is treated as its own prediction event
(no cooldown or future-dependent suppression), which reproduces the information set
a forecaster would actually face in real time.

The underlying hypothesis is that a crash reflects one of two very different things:
temporary, overreaction-driven price damage, or genuine long-term fundamental
damage. If the two are separable from contemporaneous information, then recovery
should be predictable above its base rate. We therefore frame the problem as
\emph{calibrated probability prediction} rather than as a trading system: the
priority is a defensible, leakage-free historical dataset, honest base rates, and a
measurement of whether successive information sources---price behavior, market and
sector context, point-in-time fundamentals, and later contemporaneous
filings---add predictive value. Concretely, the primary target is the probability
that a stock's close reaches at least $+10\%$ above its crash-day close at some
point within the next $20$ trading days.

\section{Data}

All market data come from Wharton Research Data Services (WRDS): CRSP for
securities, prices, and delistings, and Compustat for fundamentals. To avoid
survivorship bias, the universe is built from the full historical cross-section of
US-listed common stock---including securities that were later delisted, acquired,
renamed, or went bankrupt---rather than from today's survivors. We restrict to
common shares on NYSE, NYSE American, and Nasdaq and exclude ETFs, funds, ADRs,
warrants, and preferred shares.

We use the CRSP CIZ tables: daily prices (\texttt{dsf\_v2}), the historical security
master with full ticker history (\texttt{stocknames\_v2}), and delisting records
(\texttt{dsedelist}). Crash detection uses total returns (split- and
dividend-adjusted) so that corporate actions cannot manufacture false crashes, and
delisting returns are preserved so that terminal outcomes (e.g.\ bankruptcies) are
recorded rather than silently dropped. Fundamentals come from Compustat's unrestated
quarterly data (\texttt{comp\_urq}), dated by \emph{public availability} (the report
date, \texttt{rdq}) rather than by fiscal period end, so a quarter's figures only
become a feature after they were actually released. CRSP and Compustat are linked
through \texttt{ccmxpf\_lnkhist} (\texttt{permno}$\leftrightarrow$\texttt{gvkey}). For
the V2 extension, contemporaneous SEC filings are retrieved from EDGAR and keyed by
each filing's exact acceptance timestamp, so that only documents public by the
crash-day close are admissible.

Every feature is required to be knowable by the prediction timestamp (the crash-day
close), and provider-specific fields are normalized into canonical schemas as early
as possible so that downstream stages depend only on the schema, not on the vendor.
Table~\ref{tab:data} summarizes the resulting dataset.

\begin{table}[h]
\centering
\begin{tabular}{lr}
\toprule
Quantity & Value \\
\midrule
Securities (survivorship-safe) & 27{,}363 \\
\quad active / delisted & 4{,}171 / 23{,}192 \\
Daily price observations & $\sim$85.3 million \\
Coverage & 1925--2025 \\
Crash events (daily return $\le -10\%$) & 1{,}153{,}414 \\
\quad in liquid study pool & $\sim$268{,}000 \\
\bottomrule
\end{tabular}
\caption{Dataset summary. The liquid study pool applies point-in-time universe and
minimum-price filters computed only from information available at the event.}
\label{tab:data}
\end{table}

\end{document}
